% Part 3 导言：生成、编辑与理解
\vspace{1em}
\begin{center}
\textcolor{techgray}{\small 本部分将3DGS从重建工具扩展为内容创造和场景理解工具——物理仿真、3D生成、场景编辑、语义理解。}
\end{center}
\vspace{0.5em}

\noindent\textbf{本部分包含：}
\begin{itemize}[leftmargin=2em,itemsep=1pt]
\item[\textcolor{techorange}{\textbullet}] \textbf{第7章}（物理仿真）— 弹性体、流体、刚体、布料等多物理场仿真
\item[\textcolor{techorange}{\textbullet}] \textbf{第8章}（3D生成）— 文本到3D、图像到3D、4D动态生成
\item[\textcolor{techorange}{\textbullet}] \textbf{第9章}（场景编辑）— 语义引导编辑、风格化、几何操控
\item[\textcolor{techorange}{\textbullet}] \textbf{第10章}（场景理解）— 3D分割、开放词汇理解、目标检测、位姿估计
\end{itemize}

\noindent\textbf{阅读建议：}第8章和第9章关系密切（生成是编辑的逆操作），建议连续阅读。第7章需要第6章的动态表示基础。第10章相对独立，可随时切入。

\noindent\textbf{与其他部分的联系：}第8章的生成方法为第16章（人体化身）提供初始化方案；第10章的语义分割结果被第14章（大规模建图）和第15章（SLAM）广泛利用；第9章的编辑工具为第19章（医学影像）和第20章（遥感）提供后处理能力。

\vspace{1em}
